\documentclass[11pt,a4paper]{article}
\usepackage[utf8]{inputenc}
\usepackage[T1]{fontenc}
\usepackage{lmodern}
\usepackage{microtype}
\usepackage{geometry}
\geometry{margin=1in}
\usepackage{amsmath,amssymb,amsfonts}
\usepackage{graphicx}
\usepackage{booktabs}
\usepackage{longtable}
\usepackage{array}
\usepackage{hyperref}
\usepackage{enumitem}
\usepackage{caption}
\usepackage{xcolor}
\usepackage{listings}
\usepackage{natbib}

\hypersetup{colorlinks=true,linkcolor=blue,citecolor=blue,urlcolor=blue}

\title{Coplay-Sync: A Domain-Specific Small Language Model for Physics Research Assistance\\with LoRA, RAG, Iterative Retrieval, and Calibrated Uncertainty}
\author{Jaydip Singh \and Linkan Kumbhar}
\date{July 2026}

\begin{document}
\maketitle

\begin{abstract}
We present Coplay-Sync, a complete pipeline for building domain-specific Small Language Models (SLMs) for physics research assistance. Starting from a custom Rust-based Byte Pair Encoding (BPE) tokenizer with 32K vocabulary, we train a 35.2M parameter base model and apply Low-Rank Adaptation (LoRA) fine-tuning on 34,464 arXiv physics papers across 5 categories, achieving a 44\% evaluation loss reduction (0.0107 $\rightarrow$ 0.0060) with only 65K trainable parameters (0.19\% of base). We then integrate Retrieval-Augmented Generation (RAG) with BM25 retrieval, cross-encoder reranking, and a novel 4-component calibrated uncertainty framework that reduces mean uncertainty by 33.8\% and improves discrimination by 71.1\%. An iterative retrieval loop with 61.76\% trigger rate further refines answer quality. All components are validated through controlled ablation studies (2$\times$2 reranker $\times$ iteration design) on a 102-question physics benchmark. The complete system operates on consumer hardware (Apple M3 Pro, 16GB RAM) and is fully open-sourced.
\end{abstract}

\section{Introduction}

Large Language Models (LLMs) have demonstrated remarkable capabilities across natural language processing tasks, yet general-purpose models often lack the specialized knowledge required for physics research. Physics queries demand precise mathematical reasoning, familiarity with domain-specific notation, and access to current literature. Small Language Models (SLMs) in the 3B--7B parameter range offer a compelling alternative: they are efficient to train and deploy, yet can achieve strong performance when optimized for specific domains.

This paper introduces Coplay-Sync, a complete, reproducible pipeline for building physics research assistants. Our contributions span five phases: (1) a production Rust BPE tokenizer with 32K vocabulary, (2) base model pre-training and LoRA fine-tuning on 34K+ arXiv papers, (3) comprehensive inference and evaluation infrastructure, (4) RAG integration with reranking and calibrated uncertainty, and (5) iterative retrieval with full ablation validation.

\section{System Architecture}

The implementation follows a streaming-to-training-to-inference architecture:

\begin{enumerate}[leftmargin=*]
    \item \textbf{Data Pipeline:} Live arXiv streaming with exponential backoff retry, multi-category collection, and automatic deduplication.
    \item \textbf{Tokenizer:} Rust-based BPE with 32K vocabulary, CLI and library APIs.
    \item \textbf{Training:} Tiny Transformer with MPS/CUDA/CPU auto-selection, LoRA fine-tuning, and comprehensive checkpointing.
    \item \textbf{Inference:} Dual sampling profiles (production vs. canonical), streaming output, and multi-run evaluation.
    \item \textbf{RAG:} BM25 retrieval, cross-encoder reranking, calibrated uncertainty, and iterative refinement.
\end{enumerate}

\section{Phase 1: Tokenizer and Base Model}

\subsection{Tokenizer}

We implement byte-pair encoding (BPE) from scratch in Rust, scaling from an initial 1,000-token vocabulary to 32,000 tokens. The tokenizer handles mathematical symbols ($\alpha$, $\beta$, $\nabla$), Greek letters, subscripts/superscripts, and physics notation.

\subsection{Base Model}

The base model is a decoder-only Transformer with 35.2M parameters, 2,048 token context, and RMSNorm pre-normalization. Initial pre-training achieves an evaluation loss of 0.0107.

\section{Phase 2: LoRA Fine-Tuning}

\subsection{Training Configuration}

We apply Low-Rank Adaptation (LoRA) \citep{hu2021lora} with rank $r=8$ and $\alpha=16$ on 34,464 arXiv physics papers across 5 categories (quantum physics, optics, HEP-theory, general relativity, and general physics).

\begin{table}[h]
\centering
\caption{Phase 2 Training Specifications}
\begin{tabular}{ll}
\toprule
\textbf{Component} & \textbf{Value} \\
\midrule
Corpus Size & 34,464 documents (5 arXiv categories) \\
LoRA Rank & $r=8$, $\alpha=16$ \\
Trainable Parameters & 65,536 (0.19\% of base) \\
Training Steps & 10,000 \\
Effective Batch Size & 8 (4 micro-batch $\times$ 2 accumulation) \\
Learning Rate & 0.0002 (cosine schedule) \\
Training Time & $\sim$10.5 hours (Apple M3 Pro, MPS) \\
\bottomrule
\end{tabular}
\end{table}

\subsection{Results}

LoRA fine-tuning achieves a 44.0\% evaluation loss reduction (0.0107 $\rightarrow$ 0.0060) with 99.81\% parameter efficiency. All 11 checkpoints are ranked by validation loss; the best checkpoint at step 9,000 is selected for downstream tasks.

\begin{table}[h]
\centering
\caption{Phase 1 vs Phase 2 Performance}
\begin{tabular}{lccc}
\toprule
\textbf{Metric} & \textbf{Phase 1} & \textbf{Phase 2} & \textbf{Improvement} \\
\midrule
Eval Loss & 0.0107 & 0.0060 & \textbf{$\downarrow$44.0\%} \\
Params Trained & 35.2M & 65K & \textbf{$\downarrow$99.81\%} \\
Training Time & 48h & 10.5h & \textbf{$\downarrow$78.1\%} \\
Efficiency (loss/hour) & 0.00022 & 0.00067 & \textbf{$\uparrow$3.0$\times$} \\
\bottomrule
\end{tabular}
\end{table}

\section{Phase 3: Inference and Evaluation}

We implement dual sampling profiles for different use cases:

\begin{itemize}[leftmargin=*]
    \item \textbf{Production Profile:} temperature=2.0, top\_k=100, max\_tokens=64 (diversity-optimized)
    \item \textbf{Canonical Profile:} temperature=1.0, top\_k=50, max\_tokens=50 (reproducibility-optimized)
\end{itemize}

Both profiles show consistent LoRA acceleration ($>$1.2$\times$ speedup) with comparable inference times ($\sim$0.32--0.33s per prompt).

\section{Phase 4: RAG Integration}

\subsection{Retrieval Baseline}

BM25 retrieval achieves 100\% accuracy on the physics test corpus. We implement retry/backoff for arXiv API resilience and automatic checkpointing for long runs.

\subsection{Calibrated Uncertainty (Task 10.1)}

We introduce a 4-component calibrated uncertainty framework:

\begin{table}[h]
\centering
\caption{Calibrated Uncertainty Results}
\begin{tabular}{lcc}
\toprule
\textbf{Metric} & \textbf{Before} & \textbf{After} \\
\midrule
Mean Uncertainty & 0.662 & \textbf{0.438} ($\downarrow$33.8\%) \\
Std Dev (Discrimination) & 0.254 & \textbf{0.073} ($\downarrow$71.1\%) \\
\bottomrule
\end{tabular}
\end{table}

\subsection{Cross-Encoder Fine-Tuning (Task 10.2b)}

We collect 300 STEM preference pairs (24.3\% positive) and fine-tune a cross-encoder reranker, achieving validation loss 0.7011 and accuracy 0.7667.

\subsection{Iterative Retrieval Loop (Task 10.2c)}

The iterative loop triggers on 60\% of queries with an average of 1.6 iterations, refining answers through multi-step retrieval.

\section{Phase 5: Full Integration and Ablation}

\subsection{102-Question Benchmark}

The full integrated system is evaluated on 102 physics questions:

\begin{table}[h]
\centering
\caption{Phase 5 Full Benchmark Results}
\begin{tabular}{lc}
\toprule
\textbf{Metric} & \textbf{Value} \\
\midrule
MC Exact Match & 1.0 \\
Avg Calibrated Uncertainty & 0.4234 \\
Avg Iterations & 1.6176 \\
Trigger Rate & 61.76\% \\
Runtime (no iteration) & $\sim$25 min \\
Runtime (with iteration) & $\sim$41 min \\
\bottomrule
\end{tabular}
\end{table}

\subsection{Controlled Ablation (2$\times$2 Design)}

We isolate the pure reranker effect by disabling iteration and comparing original vs. fine-tuned rerankers. The fine-tuned reranker improves rerank scores while MC exact match remains at ceiling (1.0), indicating the need for harder benchmarks.

\section{Discussion}

The current benchmark saturates at ceiling (MC exact = 1.0) on the 102-question set, suggesting the need for harder evaluation splits. Our 2-week execution plan targets: (1) 300--500 additional hard questions, (2) multi-seed significance runs ($n=3$), (3) 1000+ strongly-labeled reranker pairs, and (4) human evaluation for faithfulness and helpfulness.

\section{Conclusion}

Coplay-Sync demonstrates that domain-specific SLMs can deliver research-grade performance at a fraction of the computational cost of general-purpose LLMs. The complete pipeline---from tokenizer to calibrated RAG---is reproducible on consumer hardware and fully open-sourced.

\section*{Acknowledgments}

We thank the open-source community for tools and datasets that enabled this work.

\begin{thebibliography}{9}
\bibitem[Hu et al.(2021)]{hu2021lora} Hu, E. J., Shen, Y., Wallis, P., et al. (2021). LoRA: Low-Rank Adaptation of Large Language Models. \textit{ICLR}.
\bibitem[Lewis et al.(2020)]{lewis2020rag} Lewis, P., Perez, E., Piktus, A., et al. (2020). Retrieval-Augmented Generation for Knowledge-Intensive NLP Tasks. \textit{NeurIPS}.
\bibitem[Touvron et al.(2023)]{touvron2023llama} Touvron, H., et al. (2023). LLaMA: Open and Efficient Foundation Language Models. \textit{arXiv:2302.13971}.
\bibitem[Vaswani et al.(2017)]{vaswani2017attention} Vaswani, A., et al. (2017). Attention Is All You Need. \textit{NeurIPS}.
\bibitem[Sennrich et al.(2016)]{sennrich2016neural} Sennrich, R., Haddow, B., Birch, A. (2016). Neural Machine Translation of Rare Words with Subword Units. \textit{ACL}.
\end{thebibliography}

\end{document}
